% ============================================================
%  DIABETES_REIMPLEMENTATION_REPORT.tex
%  Independent re-implementation of Daga et al. (2026) on the Pima Indians
%  Diabetes dataset, comparison with the published results, and re-run of
%  the repository's audit extensions on the same dataset.
%  Every table is \input from reimplementation/tex/ (generated by
%  reimplementation/compare_published.py); every figure from
%  reimplementation/figures/ (generated by reimplementation/make_figures.py).
%  Nothing numerical is typed by hand.
% ============================================================
\documentclass[11pt,a4paper]{article}
\usepackage[margin=2.1cm]{geometry}
\usepackage[T1]{fontenc}
\usepackage{lmodern}
\usepackage{amsmath,amssymb}
\usepackage{booktabs,array,longtable,tabularx,multirow}
\usepackage{graphicx}
\usepackage{xcolor}
\usepackage{caption}
\usepackage{enumitem}
\usepackage{microtype}
\usepackage[hidelinks]{hyperref}
\usepackage{float}
\usepackage{adjustbox}
\captionsetup{font=small,labelfont=bf}
\setlength{\parskip}{3pt}
\definecolor{ok}{RGB}{20,95,50}
\definecolor{warn}{RGB}{160,90,0}
\definecolor{bad}{RGB}{145,25,25}
\newcommand{\OK}[1]{\textcolor{ok}{\textbf{#1}}}
\newcommand{\WARN}[1]{\textcolor{warn}{\textbf{#1}}}
\newcommand{\BAD}[1]{\textcolor{bad}{\textbf{#1}}}
\newcommand{\code}[1]{\texttt{\small #1}}
\graphicspath{{../reimplementation/figures/}}
\newcommand{\tabdir}{../reimplementation/tex/}

\title{Independent Re-implementation of Daga et al.\ (2026) on the Pima Indians Diabetes Dataset:\\
\large Reproduction of the Published Results, Re-run of the Repository's Audit Extensions, and Independent Verification of the Findings (version 2)}
\author{Prepared for M.~Wajahat \\ \small Repository \texttt{Adaptive-Quantum-Kernel-Selection}; code in \texttt{reimplementation/}}
\date{19 October 2026 (version 2: findings edition)}

\begin{document}
\maketitle

\begin{abstract}
\noindent The pipeline of Daga, Naole, Parikh, Bommineni and Ramu, \emph{Adaptive quantum kernel selection via leakage-free stacking for clinical diagnostics on NISQ hardware} (\emph{Scientific Reports} 16:26840, 2026), was re-implemented from the paper's text for the Pima Indians Diabetes dataset ($N=768$) and run end to end: preprocessing, the three fidelity kernels, the six quantum base classifiers, the hard-voting and out-of-fold stacking ensembles, the classical baselines, the parametric and bootstrap statistics, McNemar tests, ablation, calibration, explainability, feature backtracking, the depolarizing-noise study and the variational classifier. Every published Diabetes number was then compared with the re-implementation. All 108 bootstrap values of the paper's Table~9 (nine models $\times$ four metrics $\times$ mean and both interval bounds), all 27 ablation values, the McNemar counts, the calibration metrics, the meta-learner coefficients, the SHAP, LIME and permutation rankings, the backtracked biomarker percentages and the five noise-degradation rows reproduce to the last printed digit. What does not reproduce is explained: XGBoost (GPU versus CPU build), UMAP (stochastic), kernel timing (hardware), and the variational classifier, whose released loss has an identically zero gradient so that its numbers are those of an untrained random circuit. The repository's audit extensions (ten-seed robustness, Holm correction over 32 hypotheses, fold-local preprocessing, exhaustive nested model selection, tuned classical baselines, corrected noise circuits and the gradient audit) were then applied to Diabetes, which the existing audit documents had never re-run. The recall gain of the stacking ensemble over the RBF-SVM is robust (10 of 10 seeds) and is the single quantum-superior result that survives Holm correction with $10^4$ resamples; its accuracy advantage is at parity across seeds and disappears under leakage-free, fold-local preprocessing (0.766 $\to$ 0.746 versus 0.759 for the RBF-SVM). Eight independent cross-checks (dataset identity with OpenML, circuit-level kernel checks, a verbatim run of the authors' released functions, an independent bootstrap, a code-path consistency test, two simulator checks and a hand-written Holm step-down) all pass, so the agreement with the paper is neither a transcription artefact nor a shared bug.
\end{abstract}

\tableofcontents
\clearpage

% ============================================================
\section{Findings at a glance}
\label{sec:findings}
% ============================================================
This version of the report puts the findings first; the evidence for each row is in the section named in the last column, and every number is generated by the scripts in \code{reimplementation/} (nothing is typed by hand).

\begin{table}[H]
\centering\footnotesize
\caption{Findings of the Diabetes re-implementation. ``Exact'' means identical to the paper's printed precision.}
\label{tab:findings}
\begin{adjustbox}{max width=\textwidth}
\begin{tabular}{c p{5.6cm} p{6.6cm} l}
\toprule
\# & Finding & Evidence & Where \\
\midrule
\multicolumn{4}{l}{\textit{A. What the paper reports is reproducible}} \\
1 & The headline bootstrap table (Table 9) reproduces \OK{exactly} for all nine non-variational models. & 108 of 108 values (mean and both 95\,\% bounds of accuracy, precision, recall, specificity). Stacking recall $0.621\,[0.560,0.675]$ vs RBF $0.553\,[0.493,0.616]$. & \S\ref{sec:boot} \\
2 & Ablation, McNemar, calibration, coefficients, SHAP, XAI rankings and biomarker backtracking reproduce \OK{exactly}; AUCs within 0.001. & 27/27 ablation values; $n_{01}/n_{10}=39/33$ and $35/33$; 8/8 calibration values; 8/8 coefficient/SHAP values; 12/12 ranks; 5/5 backtracked features; 15/18 AUCs exact, 3 differ by 0.001. & \S\ref{sec:auc}--\S\ref{sec:xai} \\
3 & The noise table (Table 19) reproduces \OK{exactly} under the released circuits, but its 0\,\% row is the simplified chain kernel, not the benchmark kernel. & Released: ZZ-SVM $-0.017$, stacking $+0.083$ degradation. With the benchmark kernel at 0\,\%: $+0.117$ and $+0.167$. All-to-all and chain give identical noisy accuracies. & \S\ref{sec:noise} \\
\midrule
\multicolumn{4}{l}{\textit{B. What cannot be reproduced, and why}} \\
4 & The variational classifier (VQC) is \BAD{not reproducible}: the released loss has an identically zero gradient. & $\|\nabla\|_2=0$ on both simulators; finite differences give $0.116$; 20 untrained initialisations span $0.431$--$0.650$ accuracy, containing the published $0.540$. & \S\ref{sec:vqc} \\
5 & XGBoost, UMAP and timing differ for \WARN{platform} reasons. & GPU vs CPU tree builder ($0.728$ vs $0.745$); UMAP stochastic; different simulator and CPU. & \S\ref{sec:prep}, \S\ref{sec:classical} \\
\midrule
\multicolumn{4}{l}{\textit{C. What the audit extensions add (never run on Diabetes before)}} \\
6 & The recall gain of the stacking ensemble is \OK{robust} and is the one quantum-superior result that survives family-wise correction. & Higher recall than RBF in 10 of 10 seeds (mean $0.622$ vs $0.549$). Holm $p=0.035$ with $B=10^4$; but $0.112$ with the released $B=10^3$ formula and $0.204$ with the paper's parametric tests. & \S\ref{sec:seeds}, \S\ref{sec:holm} \\
7 & The accuracy advantage is \WARN{at parity} across seeds. & RBF $0.7613\pm0.0063$ vs stacking $0.7624\pm0.0059$; RBF wins 4 seeds, stacking 6. QSVM-Angle ($0.7691$) and hard voting ($0.7676$) are higher on average than stacking. & \S\ref{sec:seeds} \\
8 & Under leakage-free fold-local preprocessing the stacking accuracy advantage \BAD{disappears}; the paper's $0.0013$ check describes the RBF baseline. & Stacking $0.766\to0.744$ (outer-local) / $0.746$ (strict); RBF $0.758\to0.759$ ($+0.0013$). 8\,\% of patients change prediction; McNemar $p=0.041$; TOST rejects $\pm1$-point equivalence. Recall gain survives ($\ge0.60$). & \S\ref{sec:foldlocal} \\
9 & QSVM-Angle, excluded from the published quartet, is the best base model in every fold, yet the fixed quartet remains the best ensemble. & Selected in 10/10 folds (accuracy, balanced accuracy); quartet chosen 0/10. Data-selected stacking $0.758$--$0.760$ vs fixed $0.766$ ($p\ge0.33$). & \S\ref{sec:nested} \\
10 & Tuning the classical baselines does not change the comparison. & Tuned RBF $0.760$, tuned XGBoost $0.759$ vs stacking $0.766$ (McNemar $p=0.72$, $0.69$). Classical SVM+KNN stacker has the highest recall of all models ($0.67$). & \S\ref{sec:fairness} \\
11 & The clinical-biomarker ranking (Table 18) does not depend on the quantum ensemble. & Identical ranking when the four meta-learner coefficients are replaced by ones. & \S\ref{sec:xai} \\
\midrule
\multicolumn{4}{l}{\textit{D. Is the agreement genuine?}} \\
12 & \OK{Yes}: eight independent cross-checks pass, including a verbatim run of the authors' released functions. & 0 prediction disagreements over 9 models; 36/36 identical CI strings; dataset identical to OpenML; kernels equal the circuits to $10^{-15}$. & \S\ref{sec:verify} \\
\bottomrule
\end{tabular}
\end{adjustbox}
\end{table}

\subsection{Bottom line}
The paper's Diabetes numbers are real and deterministic: anyone who runs the described pipeline with the stated seeds gets them to the last digit, and so does an independent implementation. The interpretation is where the picture changes. The ensemble's recall gain is genuine and robust, but it is a trade of specificity for recall produced by the balanced class weight of the meta-learner, not a sign of a better classifier: on accuracy the ensemble is tied with the untuned RBF-SVM across seeds, is beaten on average by one of its own base models (QSVM-Angle) and by simple hard voting, and loses two accuracy points once preprocessing is fitted inside the folds, while the RBF-SVM does not move. The variational-classifier results are an artefact of a zero-gradient loss, and the noise table's clean baseline is not the benchmark circuit.

\subsection{Changes from version 1}
Findings and the verification verdict moved to the front (Sections~\ref{sec:findings} and \ref{sec:verify}); the independent cross-check script \code{verify\_independent.py} and its results table were added; a per-model figure of the fold-local effect (Fig.~\ref{fig:foldlocal}) and a pipeline walkthrough (Section~\ref{sec:walkthrough}) were added; one observation outside the Diabetes scope that bears on the repository's IEEE manuscript is recorded in Appendix~\ref{app:wbco}. All numbers are unchanged from version 1 because no result file changed.

% ============================================================
\section{Is the agreement with the paper genuine? Independent cross-checks}
\label{sec:verify}
% ============================================================
An exact match between a re-implementation and a published table can arise in three unwanted ways: the published values were transcribed wrongly, the re-implementation shares a bug with the released code, or the agreement is coincidental. \code{verify\_independent.py} tests each route (Table~\ref{tab:verification}); its log is in \code{results/diabetes/verification\_report.txt}.
\begin{itemize}[nosep]
  \item \textbf{Transcription.} The published values in \code{compare\_published.py} were typed from the PDF before any result existed and were re-read against the page images of Tables 9, 10, 11, 13, 14, 16 and 19 afterwards; no discrepancy was found. Check~D confirms that the 36 published Table-9 cells equal the generated ones.
  \item \textbf{Data.} The local Pima file is row-for-row identical to OpenML's \code{diabetes} (version 1), so the input is the standard 768-row file the paper used (check~A).
  \item \textbf{Kernels.} The cached Gram matrices equal the PennyLane circuits on 600 further random pairs and on a complete row (check~B), are symmetric with unit diagonal, and are positive semi-definite.
  \item \textbf{Shared-bug risk.} The released \code{get\_oof\_predictions} and \code{run\_bootstrap} were copied verbatim (VQC lines removed) and executed on the same Gram matrices: zero prediction disagreements over the nine models and 36 of 36 identical confidence-interval strings (check~C). The re-implementation therefore computes exactly what the released code computes, and both equal the paper, which rules out a bug on either side as the source of the match.
  \item \textbf{Coincidence.} A second bootstrap with a different random-number generator and manual stratification reproduces the intervals to within 0.005 (check~E); the exact digits depend only on the seeding protocol, the conclusions do not.
  \item \textbf{Fold-local path.} The code that produced the 2-point drop under fold-local preprocessing reproduces the global predictions exactly when it is handed the globally fitted transformers (check~F), so the drop is a property of the protocol, not of the code path.
  \item \textbf{Noise simulator.} The density-matrix simulator used for the ZZ kernel equals an independent statevector evaluation of the chain circuit at $p=0$ (check~G) and PennyLane's \code{default.mixed} at $p=0.05$ for both topologies ($2\times10^{-14}$, \code{noisy\_sim.validate}).
  \item \textbf{Holm.} The adjusted $p$-values from statsmodels equal a hand-written step-down (check~H).
\end{itemize}
\input{\tabdir tab_verification.tex}
What this verification does \emph{not} cover: the GPU XGBoost builder, the stochastic UMAP embedding, kernel timing, and the paper's VQC random initialisation, which is unseeded in the released code and cannot be recovered.



% ============================================================
\section{What was done}
% ============================================================
\begin{enumerate}[nosep]
  \item \textbf{Independent re-implementation.} \code{reimplementation/qks\_common.py} implements the paper's Stage~1 preprocessing (z-score $\to$ PCA(8) $\to$ MinMax to $[0,\pi]$, fitted on the full dataset as in the released code), the three fidelity kernels of Eq.~(6) as PennyLane circuits (embedding, adjoint embedding, probability of $|0\rangle^{\otimes n}$) and as closed-form/statevector expressions, QSVM and QKNN, the hard-voting trio, the stacking quartet with its inner 5-fold out-of-fold loop and balanced logistic-regression meta-learner, the 10-fold parametric statistics, the 1,000-resample stratified bootstrap, McNemar's test and the variational classifier with Nesterov momentum. The released script \code{quantum\_clinical\_github.py} was read to settle points where the paper's text is ambiguous (Section~\ref{sec:papercode}); no code was imported from it.
  \item \textbf{Paper tables.} \code{run\_paper\_reproduction.py diabetes} regenerates Tables 4--6, 9, 10, 11, 12, 13, 14, 15, 16, 17 and 18 of the paper for Diabetes; \code{run\_noise.py diabetes} regenerates Table~19 under the released circuits and under two alternatives; \code{run\_vqc\_curves.py diabetes} regenerates Fig.~8.
  \item \textbf{Audit extensions.} \code{run\_audit\_extensions.py diabetes}, \code{run\_fold\_local\_metrics.py diabetes} and \code{run\_vqc\_gradient\_check.py diabetes} apply the repository's additional analyses (multi-seed, Holm, fold-local TOST, nested selection, tuned baselines, gradient audit) to Diabetes.
  \item \textbf{Comparison.} \code{compare\_published.py} (run with \code{QKS\_DATASETS=diabetes}) compares every number with the values transcribed from the published PDF and writes the LaTeX tables below; \code{make\_figures.py} draws the figures. Nothing in this document is typed by hand.
\end{enumerate}
The Parkinson's and WDBC pipelines were deliberately not run for this report, as requested.

\subsection{Data and environment}
Pima Indians Diabetes: 768 rows, 8 attributes, 268 diabetic (class 1) and 500 non-diabetic (class 0), downloaded from the public CSV mirror of the UCI file (the paper used the Kaggle copy of the same file; identical columns and row count). Software: Python 3.11.9, PennyLane 0.44.1 (\code{default.qubit}, \code{default.mixed}, \code{lightning.qubit}), scikit-learn 1.5.2, XGBoost 3.2.0 (CPU build), statsmodels 0.14.4, on an Intel i7-10700 (8 cores) with 16\,GB RAM. The paper used Python 3.12 on Kaggle with \code{lightning.qubit}/\code{lightning.gpu} and GPU XGBoost.

\subsection{Kernel verification}
On 100 random pairs the PennyLane circuit value $P(|0\rangle^{\otimes n})$ and the closed-form/statevector value agree to $8.9\times10^{-16}$ (Angle), $2.2\times10^{-16}$ (ZZ/IQP) and $8.9\times10^{-16}$ (Amplitude). The three $768\times768$ Gram matrices are symmetric with unit diagonal and are positive semi-definite. Explicit forms: $K_{\mathrm{Angle}}=\prod_k\cos^2\!\big((x_{ik}-x_{jk})/2\big)$; $K_{\mathrm{Amp}}=(\mathbf{x}_i^{\top}\mathbf{x}_j)^2/(\|\mathbf{x}_i\|^2\|\mathbf{x}_j\|^2)$; $K_{ZZ}=|\langle\psi(\mathbf{x}_j)|\psi(\mathbf{x}_i)\rangle|^2$ with $|\psi(\mathbf{x})\rangle=\mathrm{IQPEmbedding}(\mathbf{x})|0\rangle$ over all 28 qubit pairs.

\subsection{Pipeline walkthrough}
\label{sec:walkthrough}
The following is the exact procedure that produced every number in Section~\ref{sec:repro}; function names refer to \code{qks\_common.py}.
\begin{enumerate}[nosep]
  \item \textbf{Input.} $\mathbf{X}_{\mathrm{raw}}\in\mathbb{R}^{768\times 8}$ (Pregnancies, Glucose, BloodPressure, SkinThickness, Insulin, BMI, DiabetesPedigreeFunction, Age), $y\in\{0,1\}^{768}$ with $y=1$ for the 268 diabetic patients. Class 1 is the positive class, so \emph{recall} is the diabetic detection rate and \emph{specificity} the non-diabetic detection rate.
  \item \textbf{Stage 1 (\code{preprocess}).} $z$-score every attribute; PCA to 8 components (\code{random\_state=42}); MinMax to $[0,\pi]$. All three transformers are fitted on the full dataset before cross-validation, as in the released code (the fold-local alternative is analysed in Section~\ref{sec:foldlocal}). With 8 attributes the PCA is a rotation and keeps 100\,\% of the variance.
  \item \textbf{Stage 2 (\code{fast\_gram}, \code{circuit\_kernel}).} Three $768\times768$ Gram matrices, one per encoding, each entry $K(\mathbf{x}_i,\mathbf{x}_j)=|\langle\psi(\mathbf{x}_j)|\psi(\mathbf{x}_i)\rangle|^2$: Angle ($R_x(x_k)$ on each of 8 qubits, $K=\prod_k\cos^2((x_{ik}-x_{jk})/2)$); ZZ/IQP (Hadamards, $R_z(x_k)$, then $e^{-i x_k x_l Z_kZ_l/2}$ on all 28 pairs); Amplitude (the normalised 8-vector as the amplitudes of 3 qubits, $K$ = squared cosine similarity). Computed once for all 295,296 pairs; the circuit form (embedding, adjoint embedding, $P(|0\rangle^{\otimes n})$) agrees to $10^{-15}$.
  \item \textbf{Stage 3, base models (\code{fit\_predict}).} Outer stratified 10-fold split with shuffling, seed 42. In each fold: Classical RBF = scikit-learn \code{SVC} with the RBF kernel ($C=1$, $\gamma=$\,scale) on the scaled features; QSVM = \code{SVC} with a precomputed kernel, trained on $K[\mathrm{tr},\mathrm{tr}]$ and applied to $K[\mathrm{te},\mathrm{tr}]$; QKNN = 5-nearest-neighbours on the distance $\max(0,1-K)$.
  \item \textbf{Stage 3, ensembles (\code{inner\_oof}, \code{run\_pipeline}).} Hard voting = majority of QSVM-ZZ, QKNN-Amplitude, QSVM-Angle. OOF stacking: inside the outer training fold an inner \code{StratifiedKFold(5, seed 42)} gives each training patient an out-of-fold binary prediction from each quartet member (QSVM-Amp, QKNN-Amp, QSVM-ZZ, QKNN-Angle); a balanced logistic regression is fitted on these four binary columns and applied to the quartet's predictions on the outer test fold.
  \item \textbf{Stage 4, statistics (\code{bootstrap\_table}, \code{parametric\_table}, \code{mcnemar}).} Metrics are computed on the pooled out-of-fold predictions. Bootstrap: \code{np.random.seed(42)}, 1,000 stratified resamples of the 768 indices, 2.5/97.5 percentiles, two-sided $p=2\min\{P(\Delta\le0),P(\Delta\ge0)\}$ against RBF. Parametric: the ten fold-level metrics, $t$-interval and paired $t$-test. McNemar: exact binomial below 25 discordant pairs, continuity-corrected $\chi^2$ otherwise.
  \item \textbf{Everything else} (ablation, classical baselines, calibration, SHAP/LIME/permutation, backtracking, noise, VQC) reuses these folds and Gram matrices with the paper's stated settings; each is described where its table appears.
\end{enumerate}
\textbf{How to read the comparison tables.} Every table below has a ``published'' column transcribed from the PDF and a ``re-implementation'' column generated by the code; the ``match'' column tests equality to the paper's printed precision (three decimals, or two for Table~18). The transcriptions were checked twice against the page images.

\subsection{Where the paper's text and the released code differ}
\label{sec:papercode}
Four points had to be decided; in each case the released behaviour was reproduced because that is what produced the published tables.
\begin{enumerate}[nosep]
  \item \textbf{Angle-map rotation axis.} Eq.~(3) and Table~6 say $R_y$; the released clean kernel uses \code{qml.AngleEmbedding} with its default $R_x$. Noiseless and depolarized kernels are identical either way (Section~\ref{sec:noise}); no numerical consequence.
  \item \textbf{ZZ-map connectivity.} Eq.~(5) writes nearest-neighbour pairs; the released clean kernel uses \code{qml.IQPEmbedding}, which couples all 28 pairs. The released \emph{noise} experiment, however, implements the nearest-neighbour chain. Tables 9--18 therefore come from an all-to-all kernel and Table~19 from a chain kernel.
  \item \textbf{VQC protocol.} Fig.~8 uses an 80/20 split with 30 epochs; Tables 9 and 12 use the VQC inside the 10-fold loop with 15 in-fold epochs. Both are in the code.
  \item \textbf{Noise baseline.} The released noise code evaluates the chain circuit at every level including 0\,\%; the clean row of Table~19 is therefore not the clean benchmark kernel of Table~9.
\end{enumerate}

% ============================================================
\section{Reproduction of the published Diabetes results}
\label{sec:repro}
% ============================================================
\subsection{Preprocessing diagnostics (paper Tables 4--6)}
\label{sec:prep}
Table~\ref{tab:pca_dr_timing}: the eight components retain 100\,\% of the variance (the raw data have eight attributes), and the linear-PCA and kernel-PCA accuracies of Table~5 reproduce exactly. UMAP differs (0.728/0.738 versus 0.716/0.727); UMAP is stochastic across library versions and platforms even at a fixed seed. Kernel timing is hardware dependent and not comparable (the paper measured \code{lightning.qubit} on Kaggle; here \code{default.qubit} on a loaded desktop CPU).
\input{\tabdir tab_pca_dr_timing.tex}

\subsection{Bootstrap results (paper Table 9)}
\label{sec:boot}
Table~\ref{tab:boot_diabetes} and Fig.~\ref{fig:pub} compare the paper's 1,000-iteration bootstrap results with the re-implementation. All 108 values of the nine non-variational models (means and both 95\,\% bounds of Accuracy, Precision, Recall and Specificity) agree to three decimals, including the headline recall of the OOF stacking ensemble, $0.621\,[0.560, 0.675]$ against $0.553\,[0.493, 0.616]$ for the RBF-SVM. Because the released protocol seeds the fold split, the inner loop and the resampling, and the Gram matrices are deterministic, an independent implementation of the same pipeline reaches identical numbers. The seed-42 point estimates behind these means are: RBF-SVM 0.7578 (TN/FP/FN/TP 434/66/120/148) and OOF stacking 0.7656 (422/78/102/166).
\input{\tabdir tab_boot_diabetes.tex}
\begin{figure}[H]
\centering
\includegraphics[width=0.62\textwidth]{fig_pub_vs_reimpl.png}
\caption{Diabetes accuracy, bootstrap mean and 95\,\% interval: published Table~9 (blue circles) and re-implementation (orange squares). The two intervals coincide for every model.}
\label{fig:pub}
\end{figure}

\subsection{ROC and precision--recall areas (paper Table 10)}
\label{sec:auc}
Table~\ref{tab:auc}: 15 of 18 values agree exactly; the three remaining differ by 0.001 (RBF PR-AUC 0.703 vs 0.704, QSVM-ZZ PR-AUC 0.697 vs 0.698, QKNN-Amplitude ROC-AUC 0.752 vs 0.753), i.e.\ rounding of a quantity computed from the same predictions.
\input{\tabdir tab_auc.tex}

\subsection{Ablation (paper Table 11)}
All 27 values (nine configurations $\times$ Accuracy, Specificity, Recall) reproduce exactly (Table~\ref{tab:abl_diabetes}).
\input{\tabdir tab_abl_diabetes.tex}

\subsection{McNemar tests (paper Table 12)}
The discordant counts for stacking versus RBF ($n_{01}=39$, $n_{10}=33$, $p=0.556$) and versus the best single quantum model, QSVM-Angle (35/33, $p=0.903$ against the printed 0.904), reproduce (Table~\ref{tab:mcnemar}). The VQC row depends on an unseeded random initialisation (Section~\ref{sec:vqc}) and cannot be matched digit for digit; its magnitude (hundreds of patients the ensemble classifies correctly and the VQC does not) reproduces under both the released and the corrected loss.
\input{\tabdir tab_mcnemar.tex}

\subsection{Classical ensemble baselines (paper Table 13)}
\label{sec:classical}
Random Forest and the classical SVM+KNN stacker reproduce exactly, including the classical stacker's recall of 0.670, which is higher than the quantum ensemble's 0.621 (Table~\ref{tab:classical}). XGBoost differs (0.728 vs 0.745 accuracy): the paper used the GPU histogram implementation (\code{device="cuda"}) while this machine has no CUDA device, and the two tree builders are not bit-identical. Both values are below the quantum ensemble's 0.766.
\input{\tabdir tab_classical.tex}

\subsection{Calibration (paper Tables 14--15)}
All eight numbers (ECE and Brier, raw and after Platt scaling, for both models) reproduce exactly (Table~\ref{tab:calibration}).
\input{\tabdir tab_calibration.tex}

\subsection{Explainability and feature backtracking (paper Tables 16--18)}
\label{sec:xai}
The meta-learner coefficients, including the negative weight $-0.353$ on QKNN-Amplitude, the mean $|$SHAP$|$ values, the permutation, LIME and coefficient rankings and the top-five backtracked attributes with their percentages all reproduce exactly (Tables~\ref{tab:coef}--\ref{tab:backtrack}). The backtracking formula multiplies every PCA component's variance ratio by the same scalar $\sum_k|\beta_k|$, so the normalised percentages are independent of the meta-learner: replacing the four coefficients by ones gives the identical ranking (verified numerically). The Table~18 percentages therefore describe PCA loadings, not the quantum ensemble.
\input{\tabdir tab_coef.tex}
\input{\tabdir tab_xai.tex}
\input{\tabdir tab_backtrack.tex}

\subsection{Depolarizing-noise study (paper Table 19)}
\label{sec:noise}
Table~\ref{tab:noise} and Fig.~\ref{fig:noise}. Under the released protocol (60-patient stratified subsample, \code{default.mixed}, 5-fold CV with 3-fold inner loop, $R_y$ Angle map, 7-edge ZZ chain, Amplitude map, chain circuit also at 0\,\%) all five published degradation rows reproduce exactly, including the stacking ensemble's $+0.083/+0.125$. Two further protocols were run. (i)~\emph{Audit protocol}: the clean benchmark kernels of Table~9 at 0\,\% (this is what the repository's WDBC noise audit did). The 0\,\% QSVM-ZZ accuracy is then 0.767 instead of 0.633, so its degradation becomes $+0.117$ instead of $-0.017$ and the ensemble's $+0.167$ instead of $+0.083$: the paper's clean baseline is the simplified chain kernel, which on Diabetes is 13 points worse than the benchmark kernel before any noise is added. (ii)~\emph{Corrected circuits}: $R_x$ Angle and all-to-all ZZ under noise. The downstream accuracies at 5\,\% and 10\,\% are identical to the chain's, because the noisy QSVM-ZZ collapses to the majority class under both (accuracy 0.650, specificity 1.000 on this subsample) and the meta-learner then discards it. The $R_y$ and $R_x$ noisy Angle Gram matrices differ by $2\times10^{-16}$, so the rotation-axis discrepancy has no effect. The noisy ZZ self-fidelities differ strongly, however: at 5\,\% the chain keeps $K_{ii}\in[0.185,0.236]$ while the all-to-all circuit keeps $[0.0096,0.0153]$, close to the maximally mixed value $2^{-8}=0.0039$ (Appendix~\ref{app:noise}).
\input{\tabdir tab_noise.tex}
\begin{figure}[H]
\centering
\includegraphics[width=0.6\textwidth]{fig_noise.png}
\caption{Diabetes, 60-patient subsample, 5-fold CV: accuracy (top) and specificity (bottom) versus depolarizing probability for the five models of Table~19 under the released circuits (solid) and the benchmarked all-to-all/$R_x$ circuits (dotted). The curves coincide at 5\,\% and 10\,\%; they differ only at 0\,\%, where the released protocol uses the chain kernel.}
\label{fig:noise}
\end{figure}

\subsection{The variational classifier (paper Fig.~8, Tables 9 and 12)}
\label{sec:vqc}
The released loss builds a Python list of circuit outputs and subtracts a tensor from it; this detaches the outputs from PennyLane's autodiff graph. Table~\ref{tab:audit_vqc_grad} evaluates the gradient at the seed-42 initialisation on a fixed 16-patient Diabetes batch: the released loss returns a gradient of exactly zero on both \code{default.qubit} and \code{lightning.qubit} (the authors' device family), while central finite differences give $\|\nabla\|_2=0.116$ and the same loss with \code{pnp.stack} recovers it to $10^{-11}$. The optimizer therefore never moves the weights, and every VQC number in the paper is the output of an untrained circuit at a random, unseeded initialisation. Table~\ref{tab:vqc} shows the consequence: twenty random initialisations give 10-fold accuracies between 0.431 and 0.650 (mean 0.550), a range that contains the published 0.540; the 80/20 learning curve is flat under the released loss, and the corrected MSE and BCE losses do learn but remain below the majority rate (about 0.65) after 30 epochs (Fig.~\ref{fig:vqc}). The published ``barren plateau'' is a software defect, and the corrected-loss result is the one that should be compared with the kernel models.
\input{\tabdir tab_audit_vqc_grad.tex}
\input{\tabdir tab_vqc.tex}
\begin{figure}[H]
\centering
\includegraphics[width=0.5\textwidth]{fig_vqc_curves.png}
\caption{Diabetes VQC test accuracy over 30 epochs on the 80/20 split (paper Fig.~8 protocol): released loss (zero gradient), corrected MSE and corrected BCE. The dashed line is the majority-class rate of the test split.}
\label{fig:vqc}
\end{figure}

\subsection{Scorecard}
\begin{table}[H]
\centering\footnotesize
\caption{Reproduction scorecard for the Diabetes dataset.}
\begin{adjustbox}{max width=\textwidth}
\begin{tabular}{llp{6.6cm}}
\toprule
Paper item & Outcome & Comment \\
\midrule
Table 4 (PCA variance) & \OK{exact} & 100\,\% \\
Table 5 (DR comparison) & \OK{PCA/KPCA exact}, \WARN{UMAP differs} & UMAP is stochastic across platforms \\
Table 6 (timing) & \WARN{platform} & different simulator and CPU \\
Table 9 (bootstrap, 9 models) & \OK{108/108 exact} & means and both CI bounds \\
Table 9 (VQC row) & \BAD{not reproducible} & zero-gradient loss; RNG-dependent \\
Table 10 (AUC) & \OK{15/18 exact}, 3 within 0.001 & rounding \\
Table 11 (ablation) & \OK{27/27 exact} & \\
Table 12 (McNemar) & \OK{exact} (non-VQC) & $p$ 0.903 vs printed 0.904 \\
Table 13 (classical) & \OK{RF, classical stacker exact}; \WARN{XGBoost differs} & GPU vs CPU tree builder \\
Tables 14--15 (calibration) & \OK{8/8 exact} & \\
Table 16 (coefficients, SHAP) & \OK{8/8 exact} & \\
Table 17 (XAI rankings) & \OK{12/12 exact} & \\
Table 18 (backtracking) & \OK{5/5 exact} & ranking independent of the meta-learner \\
Table 19 (noise) & \OK{5/5 exact} under released circuits & 0\,\% row is the chain kernel \\
Fig.~8 (VQC curve) & \OK{flat curve reproduced} & because the gradient is zero \\
\bottomrule
\end{tabular}
\end{adjustbox}
\end{table}

% ============================================================
\section{The repository's audit extensions applied to Diabetes}
% ============================================================
The IEEE audit manuscript in \code{latex\_report/} (v1--v3) restricts itself to WDBC and WBCO and states that Diabetes is out of scope; the Daga-focus report only transcribes the paper's Diabetes numbers. This section applies the same analyses to Diabetes for the first time. Each analysis asks one question about the published result:
\begin{itemize}[nosep]
  \item \emph{Multi-seed robustness}: is the seed-42 ranking of the models an accident of one partition? The whole pipeline is repeated for ten outer seeds (the inner stacking seed follows the outer seed).
  \item \emph{Holm correction}: the paper reports 36 uncorrected $p$-values per dataset. Treating the 8 model $\times$ 4 metric comparisons with RBF as one family of 32 hypotheses, which survive family-wise error control at $\alpha=0.05$?
  \item \emph{Fold-local preprocessing}: the paper fits the scaler, PCA and MinMax on all patients before splitting. Does fitting them inside each training fold (so that the test fold never influences them) change the results?
  \item \emph{Nested selection}: the paper's quartet was fixed by hand. If the four base models were chosen inside each training fold by a stated criterion, which would be chosen and would it matter?
  \item \emph{Baseline fairness}: the classical baselines use default hyper-parameters. Does nested tuning change the comparison?
\end{itemize}

\subsection{Multi-seed robustness}
\label{sec:seeds}
Ten stratified 10-fold partitions (seeds 0, 1, 2, 3, 4, 5, 10, 42, 100, 999; inner CV seed equal to the outer seed) were run (Table~\ref{tab:audit_seeds}, Fig.~\ref{fig:seeds}). On accuracy the RBF-SVM and the stacking ensemble are at parity: $0.7613\pm0.0063$ versus $0.7624\pm0.0059$, 4 wins for RBF and 6 for stacking; QSVM-Angle ($0.7691$) and hard voting ($0.7676$) are both higher on average than the stacking ensemble. The headline trade-off is robust: the ensemble has the higher recall in 10 of 10 seeds (mean 0.622 versus 0.549) and the lower specificity in 10 of 10 (0.838 versus 0.875). Seed 42 is not an outlier. Using the released code's convention (inner seed fixed at 42) gives 4 wins, 1 tie and 5 losses for RBF.
\input{\tabdir tab_audit_seeds.tex}
\begin{figure}[H]
\centering
\includegraphics[width=0.55\textwidth]{fig_multiseed.png}
\caption{Diabetes: accuracy (top) and recall (bottom) of the RBF-SVM and the OOF stacking ensemble over ten cross-validation seeds.}
\label{fig:seeds}
\end{figure}

\subsection{Family-wise error control (Holm, 32 hypotheses)}
\label{sec:holm}
Eight models $\times$ four metrics against the RBF-SVM at seed 42 were tested with a stratified bootstrap of $B=10{,}000$ resamples and the $+1$-corrected two-sided $p$-value $\min\{1,2(1+\min(n_{\le},n_{\ge}))/(B+1)\}$, then Holm-adjusted (Table~\ref{tab:audit_holm}, Fig.~\ref{fig:holm}). Nine raw $p$-values are below 0.05; four survive Holm: three show QKNN-Amplitude \emph{inferior} to RBF (accuracy, precision, specificity), and one shows the stacking ensemble's recall \emph{superior} ($\Delta=+0.067$, raw $p=0.0012$, Holm $p=0.035$). This is the paper's $p=0.004$ claim (uncorrected, $B=1{,}000$), and it survives correction. The result is fragile to the testing convention: with the released formula and $B=1{,}000$ the Holm $p$ is 0.112, and with the paper's own parametric fold-level $t$-tests it is 0.204 (no survivor at all). The ensemble's accuracy, precision and specificity have Holm $p=1$.
\input{\tabdir tab_audit_holm.tex}
\begin{figure}[H]
\centering
\includegraphics[width=0.6\textwidth]{fig_holm_diabetes.png}
\caption{Diabetes: raw and Holm-adjusted bootstrap $p$-values of the 32 hypotheses versus the RBF-SVM.}
\label{fig:holm}
\end{figure}

\subsection{Global versus fold-local preprocessing}
\label{sec:foldlocal}
The paper fits the scaler, PCA and MinMax on all 768 rows before cross-validation and reports that a fold-local check changes accuracy on Diabetes by 0.0013. Two leakage-free variants were run: \emph{outer-fold-local} (transformers refitted on each outer training fold; the inner stacking loop reuses them) and \emph{strictly nested} (also refitted inside every inner fold), with the kernels recomputed from the fold-local features (Tables~\ref{tab:audit_tost} and \ref{tab:fold_local_metrics}). The RBF-SVM changes by exactly $+0.0013$ (0.7578 $\to$ 0.7591), which is the number the paper reports, but the stacking ensemble drops from 0.7656 to 0.7435 (outer-local) or 0.7461 (strict): $-2.2$ and $-2.0$ points, with 8\,\% of the patients changing prediction, McNemar $p=0.041$ and $0.078$, a paired-bootstrap 95\,\% interval of $[-0.042,-0.003]$ for the outer-local variant, and TOST $p\ge0.82$ (equivalence within $\pm1$ point is rejected). The ensemble's recall advantage over RBF survives (0.601--0.616 versus 0.549) but its accuracy advantage does not (0.744--0.746 versus 0.759); the hard-voting ensemble improves to 0.772. The published check therefore described the classical baseline, not the quantum ensemble. A mechanism is visible in the features: with fold-local MinMax the test-fold angles fall outside $[0,\pi]$ (down to $-0.70$ and up to $5.26$ rad in individual folds), which changes the fidelity kernels but hardly affects the RBF kernel.
\input{\tabdir tab_audit_tost.tex}
\input{\tabdir tab_fold_local_metrics.tex}

\begin{figure}[H]
\centering
\includegraphics[width=\textwidth]{fig_fold_local_diabetes.png}
\caption{Diabetes, seed 42: accuracy (left) and recall (right) of every model under the paper's global preprocessing and under the two fold-local variants. The RBF-SVM and the Angle and ZZ SVMs barely move; the QKNNs, the amplitude-encoded SVM and therefore the stacking ensemble do. Hard voting improves because QKNN-Amplitude, one of its voters, improves.}
\label{fig:foldlocal}
\end{figure}


\subsection{Exhaustive nested model selection}
\label{sec:nested}
Within each outer fold all $\binom{6}{4}=15$ four-model subsets were scored on the inner out-of-fold predictions (Table~\ref{tab:audit_nested}). QSVM-Angle, which the paper's quartet excludes, is selected in 10 of 10 folds by inner accuracy and by balanced accuracy (5 of 10 by meta-learner CV); the published quartet is selected in 0, 0 and 1 of 10 folds. On Diabetes the fixed quartet is nevertheless the better ensemble: 0.766 against 0.758--0.760 for the data-selected subsets (McNemar $p\ge0.33$). The selection is unstable (modal subset in 8, 9 and 2 of 10 folds; mean winning margin 0.002).
\input{\tabdir tab_audit_nested.tex}

\subsection{Classical baseline fairness}
\label{sec:fairness}
Nested grid-search tuning (inner 5-fold) raises the RBF-SVM from 0.7578 to 0.7604 and XGBoost from 0.7279 to 0.7591; neither differs significantly from the quantum ensemble's 0.7656 (McNemar $p=0.72$ and $0.69$). The classical SVM+KNN stacker keeps the highest recall of any model (0.668 released, 0.664 tuned) at lower accuracy (Table~\ref{tab:audit_fairness}).
\input{\tabdir tab_audit_fairness.tex}

% ============================================================
\section{What the repository's documents say about Diabetes}
% ============================================================
\begin{itemize}[nosep]
  \item \textbf{IEEE audit manuscript (v1--v3).} Diabetes is explicitly out of scope; none of its numbers concern this dataset. Its two code-level findings do transfer and were confirmed here: the released VQC loss has a zero gradient on Diabetes as well, and the released noise circuits differ from the benchmarked ones (with the additional observation that the paper's own 0\,\% noise row is the chain kernel).
  \item \textbf{Daga-focus report and slides.} They transcribe the paper's Diabetes values (RBF 0.758, ensemble 0.766, recall 0.553 $\to$ 0.621 at $p=0.004$, VQC 0.540, classical stacker recall 0.670). Every transcribed value is reproduced here. The report's statement that ``Angle ($R_y$)'' and ``nearest-neighbour $R_{ZZ}$'' are used follows the paper's text; the code uses $R_x$ and all-to-all coupling for the clean kernels (Section~\ref{sec:papercode}), with no numerical effect for Angle and a large effect for ZZ. Its observation that the Table~13 classical stacker beats the ensemble on recall is confirmed (0.670 versus 0.621, bootstrap means).
  \item \textbf{Comparison report v2.} Its only Diabetes-specific use of Daga et al.\ is the cross-paper comparison with Khan et al.; no Daga number is recomputed there, and none is contradicted by this re-run.
\end{itemize}

% ============================================================
\section{Conclusions}
% ============================================================
\begin{enumerate}[nosep]
  \item The paper's Diabetes results are genuine, deterministic outputs of the pipeline it describes: an independent implementation reproduces 108 of 108 bootstrap values, 27 of 27 ablation values, all McNemar counts, calibration metrics, coefficients, explainability rankings, backtracked percentages and all five noise-degradation rows to the printed precision.
  \item The exceptions are the GPU XGBoost baseline, UMAP, timing, and the variational classifier. The VQC numbers are those of an untrained circuit because the released loss has an identically zero gradient (verified on both simulator families); a corrected loss learns but stays below the majority rate within 30 epochs.
  \item The headline Diabetes claim, higher recall of the stacking ensemble, holds in 10 of 10 seeds and is the one quantum-superior result that survives Holm correction over 32 hypotheses ($B=10^4$, Holm $p=0.035$); it does not survive with the paper's parametric fold-level tests (Holm $p=0.204$).
  \item The ensemble's accuracy advantage over the RBF-SVM (0.766 versus 0.758) is at parity across seeds (4 wins, 6 losses for RBF) and is lost under leakage-free fold-local preprocessing (0.746 versus 0.759); the paper's reported 0.0013 difference corresponds to the classical baseline, not to the quantum ensemble.
  \item The published quartet omits QSVM-Angle, the best single model on Diabetes in every fold, yet remains the best stacking configuration on this dataset; data-driven selection does not improve it.
  \item Table~19 reproduces exactly, but its 0\,\% row is computed with the simplified chain kernel that the noise code implements; with the benchmark kernel at 0\,\% the ZZ-SVM degradation is $+0.117$ and the ensemble's $+0.167$. Under noise the all-to-all circuit gives the same downstream accuracies as the chain because QSVM-ZZ collapses to the majority class in both cases.
\end{enumerate}

\appendix
\section{Absolute noise results}
\label{app:noise}
\input{\tabdir tab_noise_abs.tex}


\section{An observation outside the Diabetes scope}
\label{app:wbco}
This report was restricted to Diabetes at the author's request, but one result from the WBCO run of \code{run\_audit\_extensions.py} (completed before the scope was narrowed) bears directly on the repository's IEEE manuscript (v1--v3) and is recorded here so that it is not lost; the result file is \code{audit\_H\_wbco\_rbf\_scaling.csv} in \code{results/wbco/}. The WBCO scripts in \code{wbco\_extension/} train the Classical RBF-SVM on the \emph{unscaled} PCA scores (\code{X\_wbco\_pca8.npy}), whereas the paper's pipeline, and the WDBC scripts, feed the RBF-SVM the MinMax-scaled features. With the paper's scaling, the WBCO ten-seed result changes:
\begin{center}\footnotesize
\begin{adjustbox}{max width=\textwidth}
\begin{tabular}{lccc}
\toprule
RBF input features on WBCO & RBF 10-seed mean acc. & Stacking 10-seed mean acc. & RBF wins / ties / losses \\
\midrule
Unscaled PCA scores (repository scripts) & 0.9690 & 0.9662 & 7 / 3 / 0 \\
MinMax $[0,\pi]$ scores (paper protocol) & 0.9676 & 0.9662 & 7 / 1 / 2 \\
\bottomrule
\end{tabular}
\end{adjustbox}
\end{center}
The manuscript's statement that Classical RBF ``equals or outperforms the quantum stacking ensemble in 20 out of 20 seed comparisons (zero losses)'' therefore depends on that scaling choice; under the paper's own protocol the count over WDBC and WBCO is 17 wins, 1 tie and 2 losses. The WDBC result (10 wins) is unaffected. This observation was produced by the same pipeline code that passed the Diabetes cross-checks but has not itself been cross-checked by \code{verify\_independent.py}.

\section{How to re-run}
\begin{verbatim}
cd reimplementation
python run_paper_reproduction.py diabetes
python run_noise.py diabetes 6
python run_vqc_curves.py diabetes
python run_audit_extensions.py diabetes
python run_fold_local_metrics.py diabetes
python run_vqc_gradient_check.py diabetes
python verify_independent.py
QKS_DATASETS=diabetes python compare_published.py
QKS_DATASETS=diabetes python make_figures.py
cd ../latex_report && tectonic DIABETES_REIMPLEMENTATION_REPORT_v2.tex
\end{verbatim}
Results are written to \code{reimplementation/results/diabetes/} (CSV and JSON), LaTeX fragments to \code{reimplementation/tex/} and figures to \code{reimplementation/figures/}. The lean density-matrix simulator used for the noisy ZZ kernel (\code{noisy\_sim.py}) is validated against PennyLane's \code{default.mixed} to $2\times10^{-14}$.

\end{document}
